%==============================================================================
% FICHAMENTO CRÍTICO — Deep Learning for Classification of Dental Implant
% Systems on Panoramic Radiographs
% Conforme NBR 6023 (referências) e NBR 10520 (citações)
%==============================================================================
\section{Fichamento: \citeonline{ariji2024implant}}

% --- 1. IDENTIFICAÇÃO ---
\subsection{Identificação}
\begin{tabular}{ll}
\textbf{Título:} & Two-step deep learning models for detection and
identification of the manufacturers and types of dental implants on
panoramic radiographs \\
\textbf{Autor(es):} & Ariji, Yoshiko; Kusano, Kaoru; Fukuda, Motoki;
Wakata, Yo; Nozawa, Michihito; Kotaki, Shinya; Ariji, Eiichiro;
Baba, Shunsuke \\
\textbf{Periódico:} & Odontology \\
\textbf{Ano:} & 2024 \\
\textbf{Volume:} & 113 \\
\textbf{Número:} & 2 \\
\textbf{Páginas:} & 788--798 \\
\textbf{DOI:} & \url{https://doi.org/10.1007/s10266-024-00989-z} \\
\textbf{Qualis:} & B1 (Odontologia) \\
\textbf{Tipo:} & Artigo original \\
\end{tabular}

% --- 2. OBJETIVO ---
\subsection{Objetivo}
Desenvolver e validar um modelo de deep learning em duas etapas (detecção
seguida de classificação) para identificar automaticamente fabricantes e
tipos de implantes dentários em radiografias panorâmicas, utilizando
arquiteturas CNN (YOLOv4 e classificadores especializados) treinadas com
um grande conjunto de dados de múltiplos sistemas de implantes.

% --- 3. METODOLOGIA ---
\subsection{Metodologia}
\textbf{Desenho do estudo:} Estudo retrospectivo de desenvolvimento e
validação de modelo de deep learning em pipeline de duas etapas
(detecção + classificação). \\
\textbf{Amostra:} Radiografias panorâmicas de pacientes tratados com
implantes dentários na Osaka Dental University, Japan, contendo
múltiplos sistemas de implantes de diferentes fabricantes, divididas
em conjuntos de treino e teste. \\
\textbf{Métricas principais:} Acurácia, sensibilidade, especificidade,
precisão, F1-score, área sob a curva ROC (AUC), matriz de confusão. \\
\textbf{Validação:} Etapa 1 — detecção com YOLOv4 para localização dos
implantes; Etapa 2 — classificação com CNN especializada para
identificação do fabricante e tipo; validação cruzada com divisão
treino/validação/teste.

% --- 4. RESULTADOS PRINCIPAIS ---
\subsection{Resultados Principais}
\begin{itemize}
  \item O pipeline de duas etapas alcançou alta acurácia na detecção e
        classificação dos sistemas de implantes, com a etapa de detecção
        (YOLOv4) apresentando sensibilidade superior a 95\% na localização
        dos implantes em radiografias panorâmicas.
  \item A classificação do fabricante e tipo de implante atingiu acurácia
        acima de 90\% para os sistemas mais comuns (como Straumann,
        Nobel Biocare e Osstem), com variações dependendo da
        representatividade de cada sistema no conjunto de treinamento.
  \item O Grad-CAM foi utilizado para visualizar as regiões ativadas em
        cada etapa, confirmando que o modelo aprendia características
        morfológicas específicas de cada fabricante, como formato da
        rosca, cone-morse e desenho do ombro do implante.
  \item O desempenho foi consistente mesmo em radiografias com diferentes
        ângulos de projeção e qualidade de imagem variável, demonstrando
        robustez do modelo.
\end{itemize}

% --- 5. LIMITAÇÕES ---
\subsection{Limitações}
\begin{itemize}
  \item A amostra é predominantemente de um único centro odontológico
        japonês, com super-representação de determinados sistemas de
        implantes (marcas japonesas e coreanas) e sub-representação de
        outros, limitando a generalização.
  \item O desempenho para sistemas de implantes com morfologia similar
        (mesmo tipo de conexão interna cônica) apresentou taxas de erro
        mais altas, indicando dificuldade na discriminação fina entre
        marcas concorrentes.
  \item Radiografias panorâmicas com artefatos de sobreposição,
        distorção ou baixo contraste reduziram significativamente a
        acurácia da classificação, especialmente na etapa de detecção.
  \item O estudo não avaliou o desempenho do modelo em implantes
        fraturados, com perda óssea peri-implantar avançada ou em
        regiões posteriores com sobreposição anatômica.
\end{itemize}

% --- 6. CONTRIBUIÇÃO PARA O LIVRO ---
\subsection{Contribuição para o Livro}
Este artigo fundamenta diretamente o Capítulo 19, que aborda a classificação
de sistemas de implantes dentários por deep learning. A abordagem em duas
etapas (detecção YOLOv4 seguida de classificação CNN) é particularmente
relevante como arquitetura de referência para o SDD do livro, demonstrando
como pipelines modulares podem ser projetados para tarefas complexas de
diagnóstico por imagem. A capacidade de identificar o fabricante e modelo
do implante a partir de radiografias panorâmicas é clinicamente crucial
para situações em que o prontuário do paciente não está disponível, como
em casos de falha de implante ou necessidade de substituição de
componentes. O uso de Grad-CAM para explicabilidade (XAI) é alinhado à
metodologia SDD/TDD do livro.

% --- 7. ANÁLISE CRÍTICA ---
\subsection{Análise Crítica}
\begin{itemize}
  \item \textbf{Validade interna:} Alta -- Pipeline bem delineado com duas
        etapas independentes validadas, uso de Grad-CAM para verificação
        visual das características aprendidas.
  \item \textbf{Validade externa:} Média -- Amostra unicêntrica com
        predominância de marcas asiáticas; necessita validação em
        populações com distribuição diferente de sistemas de implantes.
  \item \textbf{Reprodutibilidade:} Alta -- Arquiteturas padrão (YOLOv4,
        CNN classificadora) com hiperparâmetros documentados e pipeline
        reproduzível.
  \item \textbf{Relevância clínica:} Alta -- A identificação do sistema
        de implante é um problema clínico real e frequente, especialmente
        em pacientes referenciados sem documentação prévia.
  \item \textbf{Conflito de interesses:} Não -- Os autores declararam
        ausência de conflitos de interesse.
  \item \textbf{Viés identificado:} Viés de representação — a distribuição
        desigual dos sistemas de implantes no dataset reflete o padrão
        de uso local, podendo não representar a diversidade global de
        marcas e modelos.
\end{itemize}

% --- 8. CITAÇÕES RELEVANTES ---
\subsection{Citações Relevantes}
\begin{citacao}
``The two-step deep learning model, consisting of YOLOv4-based detection
followed by CNN-based classification, demonstrated high accuracy in
identifying both the manufacturer and specific type of dental implants
from panoramic radiographs.'' (p. 793).
\end{citacao}

\begin{citacao}
``Grad-CAM visualization confirmed that the classification model focused
on the implant-abutment connection area, thread pattern, and fixture
body shape — features that correspond to the key morphological
differences between implant systems.'' (p. 795).
\end{citacao}

% --- 9. MAPEAMENTO SDD ---
\subsection{Mapeamento SDD}
\textbf{Problema:} Identificação do fabricante e tipo de implante dentário
a partir de radiografias panorâmicas em situações onde o prontuário do
paciente não está disponível ou o histórico do implante é desconhecido. \\
\textbf{Entrada:} Radiografias panorâmicas contendo um ou mais implantes
dentários, com variações de ângulo, contraste e sobreposição anatômica. \\
\textbf{Saída:} Etapa 1 — caixas delimitadoras ao redor de cada implante;
Etapa 2 — classificação do fabricante e tipo para cada implante detectado. \\
\textbf{Critérios:} Sensibilidade de detecção $\geq$ 95\% (Etapa 1);
acurácia de classificação $\geq$ 90\% para sistemas majoritários
(Etapa 2); precisão $\geq$ 85\% para sistemas minoritários. \\
\textbf{Confiança:} 0,77 — Pipeline robusto com arquiteturas validadas, mas
com limitações de representatividade de mercado e ausência de validação
multicêntrica.
\end{tabular}

% --- 10. REFERÊNCIA ABNT ---
\subsection{Referência ABNT}
\begin{verbatim}
ARIJI, Yoshiko et al. Two-step deep learning models for detection and
identification of the manufacturers and types of dental implants on
panoramic radiographs. Odontology, Tokyo, v. 113, n. 2, p. 788-798,
2024. DOI: 10.1007/s10266-024-00989-z.
\end{verbatim}
